% ═══════════════════════════════════════════════════════════════════════
%  FICHE PEDAGOGIQUE DE MATHEMATIQUES — TERMINISH
%  Classe : 3ème — Année scolaire 2026-2027 — 8 séances de 55 min
%  Thème : CONFIGURATIONS DU PLAN
%  LEÇON 2 : PROPRIETE DE THALES
%
%  MOMENTS DIDACTIQUES APC (structure officielle, appliquée à CHAQUE
%  séance) :
%   1) Remobilisation des prérequis ou évaluation diagnostique (10 min)
%   2) Présentation de la situation (5 min)
%   3) Appropriation de la situation, compréhension de la tâche et de
%      l'organisation du travail (5 min)
%   4) Résolution du problème (individuellement 5 min puis en groupes
%      10 min)
%   5) Synthèse et bilan du travail (3 min)
%   6) Institutionnalisation : trace écrite de la leçon (7 min)
%   7) Réinvestissement / travaux dirigés (5 min)
%   8) Remédiation éventuelle (3 min)
%   9) Évaluation (formative) (2 min)
%   Total : 55 min par séance.
%
%  SITUATION D'APPRENTISSAGE : chaque question de la situation est une
%  « Synthèse » (Synthèse 1 à Synthèse 8) ; elle est le PROBLÈME résolu à
%  l'étape 4 de chaque séance. Toutes les capacités sont couvertes.
%
%  STRUCTURE DU COURS (imposée) :
%   I)   Présentation et description d'une configuration et vocabulaire
%   II)  Calculs de grandeurs
%   III) Utilisation des définitions et propriétés pour justifier
%        1) Parallélisme de deux droites (réciproque)
%   IV)  Utilisation des définitions et propriétés pour justifier
%        1) Propriété d'une configuration ; programme de construction
%
%  Compatible Overleaf (pdflatex/xelatex) ET le moteur local :
%  packages : article, geometry, xcolor, amsmath, graphicx, longtable.
% ═══════════════════════════════════════════════════════════════════════
\documentclass[10pt]{article}

\usepackage[a4paper,landscape,top=1.5cm,bottom=1.3cm,left=1.1cm,right=1.1cm]{geometry}
\usepackage{amsmath}
\usepackage[RGB]{xcolor}
\usepackage{graphicx}
\usepackage{longtable}

% ── Couleurs du canevas ────────────────────────────────────────────────
\definecolor{rougetitre}{RGB}{192,0,0}      % titres rouges soulignés
\definecolor{rougemaison}{RGB}{160,30,30}   % exercice de maison
\definecolor{vertentete}{RGB}{215,235,210}  % fond vert pâle (en-têtes)
\definecolor{bleuseance}{RGB}{46,110,180}   % cadres bleus des séances
\definecolor{bleufond}{RGB}{222,235,247}    % fond bleu pâle (n° séance)
\definecolor{jaunesurl}{RGB}{255,230,90}    % surlignage mots-clés
\definecolor{vertbilan}{RGB}{230,242,225}   % fond vert très pâle

\setlength{\parindent}{0pt}
\setlength{\tabcolsep}{3pt}
\renewcommand{\arraystretch}{1.25}

% ── Raccourcis de style ────────────────────────────────────────────────
\newcommand{\titreR}[1]{\underline{\textcolor{rougetitre}{\bfseries #1}}}
\newcommand{\app}{\titreR{Exercice d'application}}
\newcommand{\maison}{\underline{\textcolor{rougemaison}{\bfseries Exercice de maison}}}
\newcommand{\seance}[1]{\fcolorbox{bleuseance}{bleufond}{\textcolor{bleuseance}{\bfseries #1}}}
\newcommand{\moment}[1]{\textbf{\itshape #1}}
\newcommand{\mclef}[1]{\colorbox{jaunesurl}{#1}}
\newcommand{\prop}[1]{\par\medskip\noindent\fbox{\parbox{0.96\linewidth}{#1}}\par\medskip}
\newcommand{\rouge}[1]{\textcolor{rougetitre}{\bfseries #1}}
\newcommand{\ieme}{\textsuperscript{ième}}
\newcommand{\iere}{\textsuperscript{ière}}
\newcommand{\ier}{\textsuperscript{er}}
\newcommand{\ere}{\textsuperscript{ère}}

% ── En-tête, filigramme et pied de page ────────────────────────────────
\newcommand{\filigrane}{%
  \rlap{%
    \raisebox{7.5cm}[0pt][0pt]{%
      \hbox to 0pt{\hskip 3cm
        \rotatebox{30}{\scalebox{3}{\textcolor{gray!38}{\bfseries KODJONE Kodjo}}}%
      \hss}}}}
\makeatletter
\def\ps@fiche{%
  \def\@oddhead{\hfill\underline{\textcolor{rougetitre}{\bfseries FICHE PEDAGOGIQUE DE MATHEMATIQUES}}\hfill}%
  \def\@oddfoot{\filigrane
                \hfill\textcolor{bleuseance}{\bfseries M.KODJONE Kodjo Prof de Maths}\hfill
                \fcolorbox{black}{vertentete}{\arabic{page}}\hfill}%
}
\makeatother
\pagestyle{fiche}

\begin{document}

% ═════════════════════════ CARTOUCHE D'IDENTIFICATION ═════════════════════════
\begin{tabular}{|p{8.6cm}|p{8.6cm}|p{8.6cm}|}
\hline
\textbf{NOM} \; : \; {\bfseries KODJONE} &
\textbf{DRE} \; : \; {\bfseries PLO} &
\textbf{FICHE N\textsuperscript{o}} \; : \; {\bfseries 2} \\ \hline
\textbf{PRÉNOM} \; : \; {\bfseries Kodjo} &
\textbf{IESG} \; : \; {\bfseries KPALIME} &
\textbf{DISCIPLINE} \; : \; {\bfseries MATHS} \\ \hline
\textbf{CONTACT} \; : \; {\bfseries 93 150 178} &
\textbf{ÉTABLISSEMENT} \; : \; {\bfseries LYCEE KPETA} &
\textbf{CLASSE} \; : \; {\bfseries 3\ieme{}} \\ \hline
\textbf{EFFECTIF} \; : \; \textbf{G} : \dotfill \quad
                           \textbf{F} : \dotfill \quad
                           \textbf{T} : \dotfill &
\textbf{ANNÉE SCOLAIRE} \; : \; {\bfseries 2026 - 2027} &
\textbf{SEMAINE} \; : \; \dotfill \\ \hline
\end{tabular}

\bigskip
\fcolorbox{black}{vertbilan}{%
\begin{minipage}{27.1cm}
\smallskip
{\bfseries COMPETENCE 1 :} \textit{Résoudre des problèmes faisant appel aux configurations de l'espace
et du plan, aux applications du plan, à l'outil vectoriel et à la géométrie analytique.}\\[3pt]
{\bfseries THEME :} \hfill {\Large\bfseries CONFIGURATIONS DU PLAN} \hfill\mbox{}\\[3pt]
{\bfseries LEÇON 2 :} \hfill \fcolorbox{rougetitre}{white}{\textcolor{rougetitre}{\large\bfseries PROPRIETE DE THALES}} \hfill\mbox{}\\[3pt]
{\bfseries NOMBRE DE SÉANCES :} \; {\bfseries 8} \qquad
{\bfseries DURÉE D'UNE SÉANCE :} \; {\bfseries 55 min}\\[3pt]
{\bfseries SITUATION D'APPRENTISSAGE :} \; \mclef{SYNTHESES 1 à 8} (chaque question de la SA est une « Synthèse », résolue à l'étape 4 de la séance ; toutes les capacités sont couvertes)\\[3pt]
{\bfseries SUPPORTS DIDACTIQUES PRINCIPAUX :} \; Énoncé de la situation d'apprentissage ; CIAM 3\ieme{} ; CARGO 3\ieme{} ; Programme et son guide\\[3pt]
{\bfseries MATERIEL DIDACTIQUE ORDINAIRE :} \; Régie ; craie ; règle ; cahiers \qquad
{\bfseries MATERIEL DIDACTIQUE ADAPTÉ :} \; Fiche d'activités ; instruments de géométrie (règle graduée, compas, équerre)\\[3pt]
{\bfseries CONSIGNE POUR LE PROFESSEUR DE SUIVI :} \; Vérifier le respect des 9 moments didactiques APC (de la remobilisation
des prérequis à l'évaluation formative) et la trace écrite de la leçon.\\[3pt]
{\bfseries PREREQUIS :} \; Proportionnalité et quotients égaux (produit en croix) ; droites parallèles ; milieu d'un segment et
théorème des milieux (4\ieme{}) ; usage des instruments de géométrie.
\smallskip
\end{minipage}}

\bigskip
% ═════════════════════════ TABLEAU CAPACITES / CONTENUS ═════════════════════════
\begin{tabular}{|p{9.4cm}|p{17.3cm}|}
\hline
\multicolumn{1}{|c|}{\colorbox{vertentete}{\makebox[9.2cm][c]{\bfseries Capacités}}} &
\multicolumn{1}{c|}{\colorbox{vertentete}{\makebox[17.1cm][c]{\bfseries Contenus}}} \\ \hline
Reconnaître/Identifier une configuration &
\small Configurations de Thalès : \mclef{cas particulier du triangle} ; \mclef{cas général}. \hfill
\textit{(Synthèses 1 et 2)} \\ \hline
Calculer des grandeurs &
\small Longueur des segments (côtés du triangle) : \mclef{cas particulier} ; \mclef{cas général} ;
\mclef{partage d'un segment} selon une proportion donnée. \hfill \textit{(Synthèses 3, 4 et 5)} \\ \hline
Justifier &
\small \mclef{Parallélisme de deux droites} (propriété réciproque) ; propriété d'une configuration,
programme de construction. \hfill \textit{(Synthèses 6, 7 et 8)} \\ \hline
\end{tabular}

\bigskip
{\bfseries Stratégie pédagogique}

\begin{itemize}
  \item[] — \; Observation et manipulation en petits groupes
  \item[] — \; Méthode démonstrative
  \item[] — \; Travail individuel et travail en groupe
  \item[] — \; Discussions et recours aux ressources du cours (définitions, propriétés)
\end{itemize}

\bigskip
% ═════════════════════════ DEROULEMENT : TABLEAU A 4 COLONNES ═════════════════════════
\begin{longtable}{|p{2.35cm}|p{5.35cm}|p{4.55cm}|p{13.85cm}|}
% ── en-tête répété sur chaque page ──
\hline
\multicolumn{1}{|c|}{\colorbox{vertentete}{\parbox{2.15cm}{\centering\small\bfseries Moment didactique}}} &
\multicolumn{1}{c|}{\colorbox{vertentete}{\makebox[5.2cm][c]{\small\bfseries Activité de l'enseignant}}} &
\multicolumn{1}{c|}{\colorbox{vertentete}{\makebox[4.4cm][c]{\small\bfseries Activité de l'élève}}} &
\multicolumn{1}{c|}{\colorbox{vertentete}{\makebox[13.7cm][c]{\small\bfseries Trace écrite}}} \\
\hline
\endfirsthead
\hline
\multicolumn{1}{|c|}{\colorbox{vertentete}{\parbox{2.15cm}{\centering\small\bfseries Moment didactique}}} &
\multicolumn{1}{c|}{\colorbox{vertentete}{\makebox[5.2cm][c]{\small\bfseries Activité de l'enseignant}}} &
\multicolumn{1}{c|}{\colorbox{vertentete}{\makebox[4.4cm][c]{\small\bfseries Activité de l'élève}}} &
\multicolumn{1}{c|}{\colorbox{vertentete}{\makebox[13.7cm][c]{\small\bfseries Trace écrite}}} \\
\hline
\endhead
% ═══════════════════════════════════════════════════════
%  SÉANCE 1 — SYNTHESE 1 + I.1 CONFIGURATION CAS PARTICULIER
% ═══════════════════════════════════════════════════════
\seance{SEANCE 1} \newline \moment{1) Remobilisation des prérequis ou évaluation diagnostique (10 min)} &
Le professeur fait réviser la proportionnalité (quotients égaux, produit en croix) et le tracé de parallèles ; il corrige ensuite et fait le point. &
Les élèves répondent à main levée, puis certains passent au tableau. &
\titreR{Exercice prérequis} \par
\textbf{1)} Calculer $x$ : \; $\dfrac{x}{8}=\dfrac{3}{4}$. \quad
\textbf{2)} Tracer une droite $(d)$ parallèle à $(BC)$ passant par un point $M$. \par
\mclef{Rappels :} quotients égaux ; produit en croix ; équerre et règle pour tracer
une parallèle. \par
(\textit{Réponses : 1) $x=6$ ; 2) règle et équerre.}) \\ \hline

\moment{2) Présentation de la situation (5 min)} &
Le professeur recopie la situation au tableau, fait la figure, puis la lit à haute voix. &
Les élèves recopient la situation dans leur cahier et observent la figure. &
\titreR{Situation d'apprentissage} \par
{\itshape Pour aménager l'entrée du terrain de sport du LYCÉE KPETA, le maçon a tracé
deux allées rectilignes qui se coupent en $A$. Il a planté quatre piquets : $B$ et $M$
sur la première allée ($AB=9$~m ; $AM=6$~m), $C$ et $N$ sur la deuxième
($AC=7{,}5$~m ; $AN=5$~m). Il a tendu une corde entre $M$ et $N$. \textbf{Chaque
question de cette situation est appelée Synthèse (Synthèse 1, Synthèse 2, …) et sera
résolue au début du travail de chaque séance.}}
\begin{center}
\setlength{\unitlength}{1mm}%
\begin{picture}(44,32)
  \put(2,2){\rule{36mm}{0.35mm}}
  \put(20,28){\rotatebox[origin=l]{-55.30}{\rule{31.62mm}{0.35mm}}}
  \put(2,2){\rotatebox[origin=l]{55.30}{\rule{31.62mm}{0.35mm}}}
  \put(8,10.7){\rule{24mm}{0.35mm}}
  \put(19.4,27.4){\rule{1.2mm}{1.2mm}}
  \put(1.4,1.4){\rule{1.2mm}{1.2mm}}
  \put(37.4,1.4){\rule{1.2mm}{1.2mm}}
  \put(7.4,10.1){\rule{1.2mm}{1.2mm}}
  \put(31.4,10.1){\rule{1.2mm}{1.2mm}}
  \put(18.7,28.8){\small $A$}
  \put(0,-3){\small $B$}
  \put(37.5,-3){\small $C$}
  \put(3.6,11.4){\small $M$}
  \put(32.8,11.4){\small $N$}
\end{picture}
\end{center} \\ \hline

\moment{3) Appropriation de la situation, compréhension de la tâche et de
l'organisation du travail (5 min)} &
Le professeur explique les mots difficiles, pose des questions de compréhension (« Où se trouvent les points $M$ et $N$ ? Que signifie $(MN)\,//\,(BC)$ ? »), puis organise les groupes. &
Les élèves répondent aux questions, relèvent les mots difficiles et s'organisent en groupes. &
\titreR{Consignes de la Synthèse 1} \par
\textbf{1)} Repérer les points alignés et les segments $[MN]$ et $[BC]$. \par
\textbf{2)} Quelle configuration reconnaît-on ? \\ \hline

\moment{4) Résolution du problème : individuellement (5 min) puis en groupes (10 min)} &
Le professeur vérifie les productions, circule parmi les groupes et les guide sans donner la solution. &
\textbf{Résolution} : les élèves cherchent d'abord individuellement (5 min), puis confrontent leurs idées en groupes (10 min) et rédigent une production. &
\titreR{Synthèse 1 (S1)} \par
\textbf{Résolution :} \textbf{1)} $A$, $M$, $B$ sont alignés sur la première allée ;
$A$, $N$, $C$ sont alignés sur la deuxième ; $M$ et $N$ sont plus proches de $A$ que
$B$ et $C$. \par
\textbf{2)} On reconnaît une \mclef{configuration de Thalès} — cas particulier du
triangle $ABC$ : $M\in[AB]$, $N\in[AC]$. \\ \hline

\moment{5) Synthèse et bilan du travail (3 min)} &
Le professeur fait reformuler ce que la situation a permis de découvrir, corrige puis valide. &
Les élèves reformulent avec leurs propres mots, s'écoutent et se corrigent entre eux. &
\titreR{Synthèse et bilan} \par
\textbf{Données :} deux allées sécantes en $A$, piquets $B$, $M$, $C$, $N$, corde
$[MN]$. \par
\textbf{Conclusion :} la disposition des points ($M\in[AB]$, $N\in[AC]$) est celle
d'une configuration de Thalès ; il reste à connaître sa propriété. \\ \hline

\moment{6) Institutionnalisation : trace écrite de la leçon par le professeur (7 min)} &
Le professeur fait recopier la leçon (configuration et propriété, cas particulier) au tableau et dans les cahiers. &
Les élèves recopient  la trace écrite de la leçon. &
\titreR{Trace écrite} \par
\textbf{I) Présentation et description d'une configuration et vocabulaire} \par
\textbf{1) Configurations de Thalès — cas particulier du triangle} \par
— Dans le triangle $ABC$, $M$ est un point du côté $[AB]$, $N$ est un point du côté
$[AC]$ et les droites $(MN)$ et $(BC)$ sont parallèles : c'est une \mclef{configuration
de Thalès} (cas particulier du triangle).
\begin{center}
\setlength{\unitlength}{1mm}%
\begin{picture}(44,32)
  \put(2,2){\rule{36mm}{0.35mm}}
  \put(20,28){\rotatebox[origin=l]{-55.30}{\rule{31.62mm}{0.35mm}}}
  \put(2,2){\rotatebox[origin=l]{55.30}{\rule{31.62mm}{0.35mm}}}
  \put(8,10.7){\rule{24mm}{0.35mm}}
  \put(19.4,27.4){\rule{1.2mm}{1.2mm}}
  \put(1.4,1.4){\rule{1.2mm}{1.2mm}}
  \put(37.4,1.4){\rule{1.2mm}{1.2mm}}
  \put(7.4,10.1){\rule{1.2mm}{1.2mm}}
  \put(31.4,10.1){\rule{1.2mm}{1.2mm}}
  \put(18.7,28.8){\small $A$}
  \put(0,-3){\small $B$}
  \put(37.5,-3){\small $C$}
  \put(3.6,11.4){\small $M$}
  \put(32.8,11.4){\small $N$}
\end{picture}
\end{center}
\textbf{2) Propriété de Thalès (cas particulier)} \par
\prop{\textbf{\textcolor{rougetitre}{Propriété de Thalès}} \par
Si $M\in[AB]$, $N\in[AC]$ et $(MN)\,//\,(BC)$, alors \;
$\dfrac{AM}{AB}=\dfrac{AN}{AC}=\dfrac{MN}{BC}$. \par
\textit{Illustration (Synthèse 1) :} $\dfrac{AM}{AB}=\dfrac{6}{9}=\dfrac{2}{3}$ ;
$\dfrac{AN}{AC}=\dfrac{5}{7{,}5}=\dfrac{2}{3}$ : les quotients des côtés
correspondants sont égaux.} \\ \hline

\moment{7) Réinvestissement / travaux dirigés (5 min)} &
Le professeur propose l'exercice, fait passer un élève au tableau, puis corrige. &
\textbf{Résolution} : les élèves cherchent, puis deux d'entre eux passent au tableau. &
\app \par
Dans quels cas a-t-on une configuration de Thalès ? Justifier. \par
a) $M\in[AB]$, $N\in[AC]$, $(MN)\,//\,(BC)$ ; \par
b) $M\in[AB]$, $N\in[AC]$, $(MN)$ sécante avec $(BC)$. \par
\textbf{Réponses :} a) oui, cas particulier (triangle) ; b) non (pas de parallèle). \\ \hline

\moment{8) Remédiation éventuelle (3 min)} &
Le professeur reprend avec les élèves en difficulté : il fait repérer les alignements et la parallèle sur une figure simple. &
Les élèves refont  l'observation sur une nouvelle figure avec l'aide du professeur. &
\titreR{Remédiation} \par
Refaire : entourer les trois points alignés, coder la parallèle, nommer la
configuration, sur deux figures simples distribuées. \\ \hline

\moment{9) Évaluation (formative) (2 min)} &
Le professeur interroge oralement, puis donne l'exercice de maison. &
Les élèves répondent à main levée, puis notent l'exercice de maison. &
\titreR{Évaluation formative} : citer les trois éléments de la configuration de
Thalès (triangle). \par
(\textit{Réponse : $M\in[AB]$, $N\in[AC]$, $(MN)\,//\,(BC)$}.) \par
\maison \par
Tracer un triangle $ABC$ ($BC=5$~cm), placer $M\in[AB]$ et $N\in[AC]$ tels que
$(MN)\,//\,(BC)$ ; mesurer $AM$, $AB$, $AN$, $AC$, $MN$ et vérifier que
$\dfrac{AM}{AB}=\dfrac{AN}{AC}=\dfrac{MN}{BC}$. \par
\textit{(Réponse attendue : les trois quotients sont égaux.)} \\ \hline
\end{longtable}

\begin{center}
\underline{\textcolor{rougetitre}{\large\bfseries Institutionalisation \; : \; (trace écrite de la leçon par le professeur)}}
\end{center}

% ═══════════════════════════════════════════════════════
%  SÉANCE 2 — SYNTHESE 2 + I.1 CAS GÉNÉRAL + I.2 RÉCIPROQUE
% ═══════════════════════════════════════════════════════
\begin{longtable}{|p{2.35cm}|p{5.35cm}|p{4.55cm}|p{13.85cm}|}
\hline
\seance{SEANCE 2} \newline \moment{1) Remobilisation des prérequis ou évaluation diagnostique (10 min)} &
Le professeur fait corriger l'exercice de maison au tableau, puis fait le point sur les erreurs. &
\textbf{Résolution} : deux élèves présentent leur figure et leurs quotients ; les
autres corrigent. &
\titreR{Correction de l'exercice de maison} \par
Chaque élève présente sa figure et ses trois quotients ; on constate qu'ils sont
égaux (aux arrondis de mesure près). \\ \hline

\moment{2) Présentation de la situation (5 min)} &
Le professeur rappelle la situation  (figure au tableau) et lit l'énoncé de la Synthèse 2. &
Les élèves récapitulent la situation et notent l'énoncé. &
\titreR{Synthèse 2 (S2)} \par
\textbf{1)} Calculer $\dfrac{AM}{AB}$ et $\dfrac{AN}{AC}$. \par
\textbf{2)} Les points sont-ils dans le même ordre ? Que peut-on conclure pour les
droites $(MN)$ et $(BC)$ ? \\ \hline

\moment{3) Appropriation de la situation, compréhension de la tâche et de
l'organisation du travail (5 min)} &
Le professeur pose des questions de compréhension : « Comment compare-t-on deux quotients ? Que veut dire : dans le même ordre ? Que permet d'obtenir la réciproque ? » Il organise ensuite les groupes. &
Les élèves répondent, clarifient la tâche et s'organisent en groupes. &
\titreR{Consignes} : calculs séparés des deux quotients, comparaison, puis
conclusion en citant la propriété utilisée. \\ \hline

\moment{4) Résolution du problème : individuellement (5 min) puis en groupes (10 min)} &
Le professeur vérifie les productions, circule parmi les groupes et les guide sans donner la solution. &
\textbf{Résolution} : les élèves cherchent d'abord individuellement, puis en groupes ; ils rédigent une production. &
\titreR{Synthèse 2 (S2) — Résolution} \par
\textbf{1)} $\dfrac{AM}{AB}=\dfrac{6}{9}=\dfrac{2}{3}$ ; \;
$\dfrac{AN}{AC}=\dfrac{5}{7{,}5}=\dfrac{2}{3}$. \par
\textbf{2)} oui, $M$ et $N$ sont placés de $A$ vers $B$ et de $A$ vers $C$ (même
ordre) ; quotients égaux $\Rightarrow$ d'après la \mclef{réciproque de Thalès},
$(MN)\,//\,(BC)$ : la corde est parallèle au côté du terrain. \\ \hline

\moment{5) Synthèse et bilan du travail (3 min)} &
Le professeur fait reformuler la démarche, corrige puis valide. &
Les élèves reformulent  la démarche (quotients, ordre, conclusion). &
\titreR{Synthèse et bilan} : calculer les quotients, vérifier l'ordre des points,
conclure en citant la \mclef{réciproque de Thalès}. \\ \hline

\moment{6) Institutionnalisation : trace écrite de la leçon par le professeur (7 min)} &
Le professeur fait décrire la configuration générale (papillon), énonce la propriété générale puis sa réciproque, et fait recopier la leçon. &
Les élèves observent , décrivent les configurations, énoncent la propriété et recopient la
leçon. &
\titreR{Trace écrite} \par
\textbf{1) suite — Configurations de Thalès : cas général} \par
— Les points $A$, $M$, $B$ sont alignés sur une droite, les points $A$, $N$, $C$
sont alignés sur une autre droite et $(MN)\,//\,(BC)$ : configuration de Thalès
\mclef{cas général}. Dans la configuration dite du « papillon », le sommet est noté
$O$ : $O$, $A$, $B$ alignés ; $O$, $A'$, $B'$ alignés ; $(AA')\,//\,(BB')$.
\begin{center}
\setlength{\unitlength}{1mm}%
\begin{picture}(48,25)
  \put(2,2){\rotatebox[origin=l]{33.69}{\rule{28.84mm}{0.35mm}}}
  \put(2,2){\rotatebox[origin=l]{19.98}{\rule{46.82mm}{0.35mm}}}
  \put(14,10){\rule{10mm}{0.35mm}}
  \put(26,18){\rule{20mm}{0.35mm}}
  \put(1.4,1.4){\rule{1.2mm}{1.2mm}}
  \put(13.4,9.4){\rule{1.2mm}{1.2mm}}
  \put(25.4,17.4){\rule{1.2mm}{1.2mm}}
  \put(23.4,9.4){\rule{1.2mm}{1.2mm}}
  \put(45.4,17.4){\rule{1.2mm}{1.2mm}}
  \put(0,-3.4){\small $O$}
  \put(13.0,11.6){\small $A$}
  \put(23.6,6.2){\small $A'$}
  \put(24.2,19.5){\small $B$}
  \put(45.2,19.5){\small $B'$}
  \put(33,10.5){\small $(AA')\,//\,(BB')$}
\end{picture}
\end{center}
\textbf{2) Propriété de Thalès et sa réciproque} \par
\prop{\textbf{\textcolor{rougetitre}{Propriété de Thalès}} \par
Si $A$, $M$, $B$ sont alignés, $A$, $N$, $C$ sont alignés et $(MN)\,//\,(BC)$,
alors : \par
\begin{center}$\dfrac{AM}{AB}=\dfrac{AN}{AC}=\dfrac{MN}{BC}$\end{center}\par
\textbf{\textcolor{rougetitre}{Propriété réciproque}} \par
Si $A$, $M$, $B$ et $A$, $N$, $C$ sont alignés \textbf{dans le même ordre} et si
$\dfrac{AM}{AB}=\dfrac{AN}{AC}$, alors $(MN)\,//\,(BC)$.} \par
\textit{Remarque :} l'\textbf{ordre des points} est essentiel pour la réciproque. \\ \hline

\moment{7) Réinvestissement / travaux dirigés (5 min)} &
Le professeur propose l'exercice, fait passer un élève au tableau, puis corrige. &
\textbf{Résolution} : les élèves cherchent, puis deux d'entre eux passent au tableau. &
\app \par
Dans quels cas a-t-on une configuration de Thalès ? Justifier. \par
a) $O$, $A$, $B$ alignés, $O$, $A'$, $B'$ alignés, $(AA')\,//\,(BB')$ ; \par
b) $M\in[AB]$, $N\in[AC]$, $(MN)$ sécante avec $(BC)$. \par
\textbf{Réponses :} a) oui, cas général (papillon) ; b) non. \\ \hline

\moment{8) Remédiation éventuelle (3 min)} &
Le professeur reprend avec les élèves en difficulté  la notion d'ordre des points. &
Les élèves refont  un exemple simple d'alignement « dans le même ordre ». &
\titreR{Remédiation} : sur une droite graduée, placer $A$, $M$, $B$ puis $A$, $N$,
$C$ « dans le même ordre » et vérifier. \\ \hline

\moment{9) Évaluation (formative) (2 min)} &
Le professeur interroge oralement, puis donne l'exercice de maison. &
Les élèves répondent à main levée, puis notent l'exercice de maison. &
\titreR{Évaluation formative} : que permet d'obtenir la réciproque de Thalès ? \par
(\textit{Réponse : le parallélisme de deux droites}.) \par
\maison \par
\textbf{1)} Dessiner les deux configurations de Thalès (cas particulier et cas
général) en codant la parallèle. \par
\textbf{2)} Réciter la propriété de Thalès et sa réciproque. \par
\textit{(Réponse attendue : deux figures correctes ; énoncés conformes.)} \\ \hline
\end{longtable}

\begin{center}
\underline{\textcolor{rougetitre}{\large\bfseries Institutionalisation \; : \; (trace écrite de la leçon par le professeur)}}
\end{center}

% ═══════════════════════════════════════════════════════
%  SÉANCE 3 — SYNTHESE 3 + II.1 CALCULS : CAS PARTICULIER
% ═══════════════════════════════════════════════════════
\begin{longtable}{|p{2.35cm}|p{5.35cm}|p{4.55cm}|p{13.85cm}|}
\hline
\seance{SEANCE 3} \newline \moment{1) Remobilisation des prérequis ou évaluation diagnostique (10 min)} &
Le professeur fait corriger l'exercice de maison au tableau, puis fait le point sur les erreurs. &
\textbf{Résolution} : deux élèves présentent leurs figures et les énoncés ; les
autres corrigent. &
\titreR{Correction de l'exercice de maison} \par
On vérifie les deux figures (parallèle codée) et les deux énoncés (propriété et réciproque). \\ \hline

\moment{2) Présentation de la situation (5 min)} &
Le professeur rappelle la situation  (figure au tableau) et lit l'énoncé de la Synthèse 3. &
Les élèves récapitulent la situation et notent l'énoncé. &
\titreR{Synthèse 3 (S3)} : la corde mesure $MN=4$~m. Calculer $BC$. \\ \hline

\moment{3) Appropriation de la situation, compréhension de la tâche et de
l'organisation du travail (5 min)} &
Le professeur pose des questions de compréhension : « Quelle propriété peut-on utiliser ? Quel quotient est déjà connu ? » Il organise ensuite les groupes. &
Les élèves répondent, clarifient la tâche et s'organisent en groupes. &
\titreR{Consignes} : écrire la propriété de Thalès, remplacer les longueurs
connues, conclure par un produit en croix. \\ \hline

\moment{4) Résolution du problème : individuellement (5 min) puis en groupes (10 min)} &
Le professeur vérifie les productions, circule parmi les groupes et les guide sans donner la solution. &
\textbf{Résolution} : les élèves cherchent d'abord individuellement, puis en groupes ; ils rédigent une production. &
\titreR{Synthèse 3 (S3) — Résolution} \par
$(MN)\,//\,(BC)$ (établie à la Synthèse 2), donc d'après la propriété de Thalès : \;
$\dfrac{MN}{BC}=\dfrac{AM}{AB}=\dfrac{2}{3}$. \par
$BC=\dfrac{MN\times3}{2}=\dfrac{4\times3}{2}=6$~m. \par
Le côté du terrain mesure $6$~m. \\ \hline

\moment{5) Synthèse et bilan du travail (3 min)} &
Le professeur fait reformuler la démarche, corrige puis valide. &
Les élèves reformulent  la démarche (propriété, produit en croix). &
\titreR{Synthèse et bilan} : la propriété de Thalès permet de calculer une longueur
inconnue par un produit en croix. \\ \hline

\moment{6) Institutionnalisation : trace écrite de la leçon par le professeur (7 min)} &
Le professeur présente la méthode de calcul des longueurs dans le cas particulier, fait des exemples rédigés au tableau, puis fait recopier la leçon. &
Les élèves observent  la méthode, calculent sous la dictée et recopient la leçon. &
\titreR{Trace écrite} \par
\textbf{II) Calculs de grandeurs} \par
\textbf{1) Longueur des segments, côtés du triangle (propriété de Thalès) — cas
particulier} \par
— \mclef{Méthode :} \textbf{1)} identifier la configuration ($M\in[AB]$,
$N\in[AC]$, $(MN)\,//\,(BC)$) ; \textbf{2)} écrire la propriété de Thalès ;
\textbf{3)} remplacer les longueurs connues ; \textbf{4)} conclure par un produit en
croix. \par
\textit{Exemple 1.} $AM=3$~cm ; $AB=7{,}5$~cm ; $AN=2{,}4$~cm ; $BC=5$~cm. \par
$\dfrac{AM}{AB}=\dfrac{MN}{BC}$ : $MN=\dfrac{AM\times BC}{AB}
=\dfrac{3\times5}{7{,}5}=2$~cm. \par
$\dfrac{AM}{AB}=\dfrac{AN}{AC}$ : $AC=\dfrac{AB\times AN}{AM}
=\dfrac{7{,}5\times2{,}4}{3}=6$~cm. \par
\textit{Exemple 2.} Calculer $AB$ sachant $AM=4$~cm, $AN=3$~cm, $AC=9$~cm : \;
$AB=\dfrac{AC\times AM}{AN}=\dfrac{9\times4}{3}=12$~cm. \\ \hline

\moment{7) Réinvestissement / travaux dirigés (5 min)} &
Le professeur propose l'exercice, fait passer un élève au tableau, puis corrige. &
\textbf{Résolution} : les élèves cherchent, puis deux d'entre eux passent au tableau. &
\app \par
$ABC$ est un triangle, $M\in[AB]$, $N\in[AC]$, $(MN)\,//\,(BC)$. \par
\textbf{1)} $AB=6$~cm ; $AM=4$~cm ; $AC=9$~cm ; $BC=7{,}5$~cm. Calculer $AN$ et $MN$. \par
\textbf{2)} $AM=5$~cm ; $AB=12$~cm ; $AN=3{,}5$~cm. Calculer $AC$. \par
\textbf{Réponses :} \textbf{1)} $AN=\dfrac{9\times4}{6}=6$~cm ; \;
$MN=\dfrac{7{,}5\times4}{6}=5$~cm. \par
\textbf{2)} $AC=\dfrac{12\times3{,}5}{5}=8{,}4$~cm. \\ \hline

\moment{8) Remédiation éventuelle (3 min)} &
Le professeur reprend avec les élèves en difficulté  le produit en croix. &
Les élèves refont  un produit en croix simple. &
\titreR{Remédiation} : calculer $x$ sachant $\dfrac{x}{10}=\dfrac{3}{5}$ ;
(\textit{réponse : $x=6$}). \\ \hline

\moment{9) Évaluation (formative) (2 min)} &
Le professeur interroge oralement, puis donne l'exercice de maison. &
Les élèves répondent à main levée, puis notent l'exercice de maison. &
\titreR{Évaluation formative} : quel outil utilise-t-on pour calculer une longueur
dans une configuration de Thalès ? \par
(\textit{Réponse : la propriété de Thalès, puis un produit en croix}.) \par
\maison \par
$M\in[AB]$, $N\in[AC]$, $(MN)\,//\,(BC)$ ; $AM=2{,}5$~cm ; $AB=10$~cm ;
$AC=12$~cm ; $MN=2$~cm. \par
\textbf{1)} Calculer $AN$. \quad \textbf{2)} Calculer $BC$. \par
\textit{(Réponse attendue : $AN=3$~cm et $BC=8$~cm.)} \\ \hline
\end{longtable}

\begin{center}
\underline{\textcolor{rougetitre}{\large\bfseries Institutionalisation \; : \; (trace écrite de la leçon par le professeur)}}
\end{center}

% ═══════════════════════════════════════════════════════
%  SÉANCE 4 — SYNTHESE 4 + II.1 CALCULS : CAS GÉNÉRAL
% ═══════════════════════════════════════════════════════
\begin{longtable}{|p{2.35cm}|p{5.35cm}|p{4.55cm}|p{13.85cm}|}
\hline
\seance{SEANCE 4} \newline \moment{1) Remobilisation des prérequis ou évaluation diagnostique (10 min)} &
Le professeur fait corriger l'exercice de maison au tableau, puis fait le point sur les erreurs. &
\textbf{Résolution} : deux élèves traitent l'exercice au tableau ; les autres corrigent. &
\titreR{Correction de l'exercice de maison} \par
\textbf{1)} $AN=\dfrac{AC\times AM}{AB}=\dfrac{12\times2{,}5}{10}=3$~cm. \par
\textbf{2)} $BC=\dfrac{AB\times MN}{AM}=\dfrac{10\times2}{2{,}5}=8$~cm. \\ \hline

\moment{2) Présentation de la situation (5 min)} &
Présente la deuxième partie de la situation (papillon) et lit l'énoncé de la
Synthèse 4. &
Les élèves recopient  l'énoncé et observent le schéma. &
\titreR{Synthèse 4 (S4)} \par
{\itshape Deuxième partie : pour l'autre accès du terrain, deux chemins rectilignes
se coupent en $O$. Des piquets $A$, $B$ sont plantés sur le premier chemin
($OA=4$~m ; $OB=10$~m) et des piquets $A'$, $B'$ sur le deuxième ($OA'=6$~m).
On sait que $(AA')\,//\,(BB')$.} \par
\textbf{1)} Quelle configuration reconnaît-on ? \textbf{2)} Calculer $OB'$. \\ \hline

\moment{3) Appropriation de la situation, compréhension de la tâche et de
l'organisation du travail (5 min)} &
Le professeur pose des questions de compréhension : « Quelles droites sont sécantes ? Quelle formule doit-on écrire ? » Il organise ensuite les groupes. &
Les élèves répondent, clarifient la tâche et s'organisent en groupes. &
\titreR{Consignes} : identifier les droites sécantes en $O$ et la droite parallèle,
écrire les quotients égaux, conclure. \\ \hline

\moment{4) Résolution du problème : individuellement (5 min) puis en groupes (10 min)} &
Le professeur vérifie les productions, circule parmi les groupes et les guide sans donner la solution. &
\textbf{Résolution} : les élèves cherchent d'abord individuellement, puis en groupes ; ils rédigent une production. &
\titreR{Synthèse 4 (S4) — Résolution} \par
\textbf{1)} configuration de Thalès \textbf{cas général} (papillon) : droites
sécantes en $O$ et droite parallèle. \par
\textbf{2)} $\dfrac{OA}{OB}=\dfrac{OA'}{OB'}$ :
$OB'=\dfrac{OB\times OA'}{OA}=\dfrac{10\times6}{4}=15$~m. \\ \hline

\moment{5) Synthèse et bilan du travail (3 min)} &
Le professeur fait reformuler la démarche, corrige puis valide. &
Les élèves reformulent  la démarche (configuration, quotients, produit en croix). &
\titreR{Synthèse et bilan} : dans le cas général, on écrit
$\dfrac{OA}{OB}=\dfrac{OA'}{OB'}$ puis on fait un produit en croix. \\ \hline

\moment{6) Institutionnalisation : trace écrite de la leçon par le professeur (7 min)} &
Le professeur présente le calcul des longueurs dans la configuration générale (papillon), fait des exemples rédigés avec la figure, puis fait recopier la leçon. &
Les élèves observent  la figure, calculent sous la dictée et recopient la leçon. &
\titreR{Trace écrite} \par
\textbf{1) suite — Longueur des segments, côtés du triangle : cas général} \par
— Dans la configuration générale, on utilise les droites sécantes en $O$ : \;
$\dfrac{OA}{OB}=\dfrac{OA'}{OB'}=\dfrac{AA'}{BB'}$. \par
\begin{center}
\setlength{\unitlength}{1mm}%
\begin{picture}(48,25)
  \put(2,2){\rotatebox[origin=l]{33.69}{\rule{28.84mm}{0.35mm}}}
  \put(2,2){\rotatebox[origin=l]{19.98}{\rule{46.82mm}{0.35mm}}}
  \put(14,10){\rule{10mm}{0.35mm}}
  \put(26,18){\rule{20mm}{0.35mm}}
  \put(1.4,1.4){\rule{1.2mm}{1.2mm}}
  \put(13.4,9.4){\rule{1.2mm}{1.2mm}}
  \put(25.4,17.4){\rule{1.2mm}{1.2mm}}
  \put(23.4,9.4){\rule{1.2mm}{1.2mm}}
  \put(45.4,17.4){\rule{1.2mm}{1.2mm}}
  \put(0,-3.4){\small $O$}
  \put(13.0,11.6){\small $A$}
  \put(23.6,6.2){\small $A'$}
  \put(24.2,19.5){\small $B$}
  \put(45.2,19.5){\small $B'$}
\end{picture}
\end{center}
\textit{Exemple 1.} $(AA')\,//\,(BB')$ ; $OA=4$~cm ; $OB=10$~cm ; $OA'=6$~cm. \par
$\dfrac{OA}{OB}=\dfrac{OA'}{OB'}$ : $OB'=\dfrac{OB\times OA'}{OA}
=\dfrac{10\times6}{4}=15$~cm. \par
\textit{Exemple 2.} Si de plus $BB'=12$~cm : \;
$AA'=\dfrac{OA\times BB'}{OB}=\dfrac{4\times12}{10}=4{,}8$~cm. \\ \hline

\moment{7) Réinvestissement / travaux dirigés (5 min)} &
Le professeur propose l'exercice, fait passer un élève au tableau, puis corrige. &
\textbf{Résolution} : les élèves cherchent, puis deux d'entre eux passent au tableau. &
\app \par
$O$, $A$, $B$ alignés ; $O$, $A'$, $B'$ alignés ; $(AA')\,//\,(BB')$. \par
\textbf{1)} $OA=3$~cm ; $OB=7{,}5$~cm ; $OA'=4{,}5$~cm. Calculer $OB'$. \par
\textbf{Réponse :} $OB'=\dfrac{7{,}5\times4{,}5}{3}=11{,}25$~cm. \\ \hline

\moment{8) Remédiation éventuelle (3 min)} &
Le professeur reprend avec les élèves en difficulté  l'écriture des quotients dans le papillon. &
Les élèves refont  l'écriture des quotients sur un schéma simple. &
\titreR{Remédiation} : compléter $\dfrac{OA}{OB}=\dfrac{\dots}{\dots}$ sur un
schéma, puis calculer la longueur manquante. \\ \hline

\moment{9) Évaluation (formative) (2 min)} &
Le professeur interroge oralement, puis donne l'exercice de maison. &
Les élèves répondent à main levée, puis notent l'exercice de maison. &
\titreR{Évaluation formative} : écrire la relation des quotients dans la
configuration générale. \par
(\textit{Réponse : $\dfrac{OA}{OB}=\dfrac{OA'}{OB'}=\dfrac{AA'}{BB'}$}.) \par
\maison \par
Configuration générale avec $(AA')\,//\,(BB')$ : $OA=5$~cm ; $OB=9$~cm ;
$OB'=18$~cm ; $BB'=7{,}2$~cm. \par
\textbf{1)} Calculer $OA'$. \quad \textbf{2)} Calculer $AA'$. \par
\textit{(Réponse attendue : $OA'=10$~cm et $AA'=4$~cm.)} \\ \hline
\end{longtable}

\begin{center}
\underline{\textcolor{rougetitre}{\large\bfseries Institutionalisation \; : \; (trace écrite de la leçon par le professeur)}}
\end{center}

% ═══════════════════════════════════════════════════════
%  SÉANCE 5 — SYNTHESE 5 + II.2 PARTAGE D'UN SEGMENT
% ═══════════════════════════════════════════════════════
\begin{longtable}{|p{2.35cm}|p{5.35cm}|p{4.55cm}|p{13.85cm}|}
\hline
\seance{SEANCE 5} \newline \moment{1) Remobilisation des prérequis ou évaluation diagnostique (10 min)} &
Le professeur fait corriger l'exercice de maison au tableau, puis fait le point sur les erreurs. &
\textbf{Résolution} : deux élèves traitent l'exercice au tableau ; les autres corrigent. &
\titreR{Correction de l'exercice de maison} \par
\textbf{1)} $OA'=\dfrac{OA\times OB'}{OB}=\dfrac{5\times18}{9}=10$~cm. \par
\textbf{2)} $AA'=\dfrac{5\times7{,}2}{9}=4$~cm. \\ \hline

\moment{2) Présentation de la situation (5 min)} &
Le professeur rappelle la situation  et lit l'énoncé de la Synthèse 5. &
Les élèves récapitulent la situation et notent l'énoncé. &
\titreR{Synthèse 5 (S5)} : le maçon veut placer la porte $P$ du terrain en
partageant le côté $[BC]$ ($BC=6$~m) dans la proportion $1:2$. \par
\textbf{1)} Calculer les deux parts. \textbf{2)} Écrire le programme de
construction. \\ \hline

\moment{3) Appropriation de la situation, compréhension de la tâche et de
l'organisation du travail (5 min)} &
Le professeur pose des questions de compréhension : « Que signifie partager un segment dans la proportion $1:2$ ? Quels instruments utiliser ? » Il organise ensuite les groupes. &
Les élèves répondent, clarifient la tâche et s'organisent en groupes. &
\titreR{Consignes} : calculer les parts ($\dfrac{1}{3}$ et $\dfrac{2}{3}$ de $BC$),
puis décrire la construction (demi-droite, reports égaux, parallèle). \\ \hline

\moment{4) Résolution du problème : individuellement (5 min) puis en groupes (10 min)} &
Le professeur vérifie les productions, circule parmi les groupes et les guide sans donner la solution. &
\textbf{Résolution} : les élèves calculent et construisent d'abord individuellement, puis en groupes ; ils rédigent une production. &
\titreR{Synthèse 5 (S5) — Résolution} \par
\textbf{1)} $BP=\dfrac{1}{3}\times6=2$~m ; $PC=\dfrac{2}{3}\times6=4$~m
($2+4=6$~m). \par
\textbf{2)} tracer une demi-droite $[Bx)$ ; reporter $3$ longueurs égales ($1+2$) ;
relier le $3$\ieme{} point à $C$ ; tracer la parallèle passant par le $1$\ier{}
point : elle coupe $[BC]$ en $P$. \\ \hline

\moment{5) Synthèse et bilan du travail (3 min)} &
Le professeur fait reformuler la démarche, corrige puis valide. &
Les élèves reformulent  la démarche (parts, programme de construction). &
\titreR{Synthèse et bilan} : partager un segment, c'est calculer des fractions de la
longueur totale ou construire par reports égaux (propriété de Thalès). \\ \hline

\moment{6) Institutionnalisation : trace écrite de la leçon par le professeur (7 min)} &
Le professeur présente le partage d'un segment selon une proportion, fait construire la figure, énonce le programme de construction, puis fait recopier la leçon. &
Les élèves calculent, construisent la figure pas à pas, puis recopient la leçon. &
\titreR{Trace écrite} \par
\textbf{2) Partage d'un segment selon une proportion donnée (propriété de Thalès)} \par
— \mclef{Partager} un segment $[AB]$ dans la proportion $m:n$, c'est placer un point
$M\in[AB]$ tel que $\dfrac{AM}{MB}=\dfrac{m}{n}$, c'est-à-dire $AM=\dfrac{m}{m+n}\times AB$. \par
\textit{Exemple.} Partager $[AB]$ de longueur $10$~cm dans la proportion $2:3$ : \;
$AM=\dfrac{2}{5}\times10=4$~cm et $MB=6$~cm. \par
— \mclef{Programme de construction} (sans calculer) : \par
\textbf{1)} tracer une demi-droite $[Ax)$ ; \par
\textbf{2)} reporter sur $[Ax)$, à partir de $A$, $5$ longueurs égales
($2+3$) : $M_1$, $M_2$, …, $M_5$ ; \par
\textbf{3)} relier $M_5$ à $B$ ; \par
\textbf{4)} tracer la parallèle à $(M_5B)$ passant par $M_2$ : elle coupe $[AB]$ en
$M$ ; \par
\textbf{5)} on a $\dfrac{AM}{AB}=\dfrac{AM_2}{AM_5}=\dfrac{2}{5}$ (propriété de
Thalès), donc $AM=4$~cm.
\begin{center}
\setlength{\unitlength}{1mm}%
\begin{picture}(56,36)
  \put(2,2){\rule{50mm}{0.35mm}}
  \put(2,2){\rotatebox[origin=l]{45}{\rule{42.4mm}{0.35mm}}}
  \put(32,32){\rotatebox[origin=l]{-56.31}{\rule{36.06mm}{0.35mm}}}
  \put(14,14){\rotatebox[origin=l]{-56.31}{\rule{14.42mm}{0.35mm}}}
  \put(7.4,7.4){\rule{1.2mm}{1.2mm}}
  \put(13.4,13.4){\rule{1.2mm}{1.2mm}}
  \put(19.4,19.4){\rule{1.2mm}{1.2mm}}
  \put(25.4,25.4){\rule{1.2mm}{1.2mm}}
  \put(31.4,31.4){\rule{1.2mm}{1.2mm}}
  \put(21.4,1.4){\rule{1.2mm}{1.2mm}}
  \put(0,-3.4){\small $A$}
  \put(51,-3.4){\small $B$}
  \put(20.6,4.8){\small $M$}
  \put(9.2,8.6){\scriptsize $M_1$}
  \put(15.2,14.6){\scriptsize $M_2$}
  \put(21.2,20.6){\scriptsize $M_3$}
  \put(27.2,26.6){\scriptsize $M_4$}
  \put(33.2,32.4){\scriptsize $M_5$}
\end{picture}
\end{center} \\ \hline

\moment{7) Réinvestissement / travaux dirigés (5 min)} &
Le professeur propose l'exercice, fait travailler les élèves individuellement, puis corrige. &
\textbf{Résolution} : les élèves calculent et construisent, puis passent au tableau. &
\app \par
\textbf{1)} Calculer les parts d'un segment $[AB]$ de $9$~cm partagé dans la
proportion $3:2$. \par
\textbf{Réponse :} $AM=\dfrac{3}{5}\times9=5{,}4$~cm et
$MB=\dfrac{2}{5}\times9=3{,}6$~cm ($5{,}4+3{,}6=9$~cm). \\ \hline

\moment{8) Remédiation éventuelle (3 min)} &
Le professeur reprend avec les élèves en difficulté  le calcul d'une fraction d'une longueur. &
Les élèves refont  un calcul de part. &
\titreR{Remédiation} : calculer $\dfrac{2}{7}$ de $21$~cm ;
(\textit{réponse : $6$~cm}). \\ \hline

\moment{9) Évaluation (formative) (2 min)} &
Le professeur interroge oralement, puis donne l'exercice de maison. &
Les élèves répondent à main levée, puis notent l'exercice de maison. &
\titreR{Évaluation formative} : citer les étapes du programme de construction. \par
(\textit{Réponse : demi-droite ; reports égaux ; droite ; parallèle}.) \par
\maison \par
\textbf{1)} Partager un segment $[AB]$ de $8$~cm dans la proportion $1:3$
(calcul des parts). \par
\textbf{2)} Écrire le programme de construction correspondant. \par
\textit{(Réponse attendue : $AM=2$~cm et $MB=6$~cm ; programme en $4$ reports égaux, avec la parallèle passant par $M_1$.)} \\ \hline
\end{longtable}

\begin{center}
\underline{\textcolor{rougetitre}{\large\bfseries Institutionalisation \; : \; (trace écrite de la leçon par le professeur)}}
\end{center}

% ═══════════════════════════════════════════════════════
%  SÉANCE 6 — SYNTHESE 6 + III) PARALLÉLISME (RÉCIPROQUE)
% ═══════════════════════════════════════════════════════
\begin{longtable}{|p{2.35cm}|p{5.35cm}|p{4.55cm}|p{13.85cm}|}
\hline
\seance{SEANCE 6} \newline \moment{1) Remobilisation des prérequis ou évaluation diagnostique (10 min)} &
Le professeur fait corriger l'exercice de maison au tableau, puis fait le point sur les erreurs. &
\textbf{Résolution} : deux élèves traitent l'exercice au tableau ; les autres corrigent. &
\titreR{Correction de l'exercice de maison} \par
\textbf{1)} $AM=\dfrac{1}{4}\times8=2$~cm ; $MB=\dfrac{3}{4}\times8=6$~cm. \par
\textbf{2)} demi-droite $[Ax)$ ; $4$ reports égaux ; relier $M_4$ à $B$ ; parallèle
à $(M_4B)$ passant par $M_1$. \\ \hline

\moment{2) Présentation de la situation (5 min)} &
Le professeur rappelle la situation  et lit l'énoncé de la Synthèse 6. &
Les élèves récapitulent la situation et notent l'énoncé. &
\titreR{Synthèse 6 (S6)} \par
{\itshape Sur le second plan du terrain : $OA=4$~m ; $OB=8$~m ; $OA'=6$~m ;
$OB'=12$~m. Le maçon affirme que les cheminements $(AA')$ et $(BB')$ sont
parallèles. A-t-il raison ? Justifier.} \\ \hline

\moment{3) Appropriation de la situation, compréhension de la tâche et de
l'organisation du travail (5 min)} &
Le professeur pose des questions de compréhension : « Que faut-il vérifier avant d'utiliser la réciproque ? » Il organise ensuite les groupes. &
Les élèves répondent, clarifient la tâche et s'organisent en groupes. &
\titreR{Consignes} : vérifier l'ordre des points, calculer séparément les deux
quotients, comparer, conclure en citant la réciproque. \\ \hline

\moment{4) Résolution du problème : individuellement (5 min) puis en groupes (10 min)} &
Le professeur vérifie les productions, circule parmi les groupes et les guide sans donner la solution. &
\textbf{Résolution} : les élèves cherchent d'abord individuellement, puis en groupes ; ils rédigent une production. &
\titreR{Synthèse 6 (S6) — Résolution} \par
Les points sont alignés dans le même ordre ; \par
$\dfrac{OA}{OB}=\dfrac{4}{8}=0{,}5$ et $\dfrac{OA'}{OB'}=\dfrac{6}{12}=0{,}5$. \par
Quotients égaux $\Rightarrow$ d'après la \mclef{réciproque de Thalès},
$(AA')\,//\,(BB')$ : le maçon a raison. \\ \hline

\moment{5) Synthèse et bilan du travail (3 min)} &
Le professeur fait reformuler la démarche de justification, corrige puis valide. &
Les élèves reformulent  les quatre étapes de la justification. &
\titreR{Synthèse et bilan} : la réciproque de Thalès sert à \mclef{justifier}
(prouver) que deux droites sont parallèles. \\ \hline

\moment{6) Institutionnalisation : trace écrite de la leçon par le professeur (7 min)} &
Le professeur présente la méthode de justification du parallélisme à l'aide de la réciproque ; il traite au tableau un exemple rédigé et un contre-exemple, puis fait recopier la leçon. &
Les élèves observent  la méthode, rédigent sous la dictée et recopient la leçon. &
\titreR{Trace écrite} \par
\textbf{III) Utilisation des définitions et propriétés pour justifier} \par
\textbf{1) Parallélisme de deux droites (propriété réciproque)} \par
— \mclef{Méthode :} \textbf{1)} vérifier que les points sont alignés \textbf{dans le
même ordre} ; \textbf{2)} calculer séparément $\dfrac{AM}{AB}$ et $\dfrac{AN}{AC}$ ;
\textbf{3)} comparer les deux quotients ; \textbf{4)} conclure en citant la
réciproque de Thalès. \par
\textit{Exemple 1.} $OA=4$~cm ; $OB=8$~cm ; $OA'=6$~cm ; $OB'=12$~cm. \par
$\dfrac{OA}{OB}=0{,}5$ ; \; $\dfrac{OA'}{OB'}=0{,}5$ ;
points dans le même ordre : d'après la \mclef{réciproque de Thalès},
$(AA')\,//\,(BB')$. \par
\textit{Exemple 2 (contre-exemple).} $OA=4$ ; $OB=10$ ; $OA'=6$ ; $OB'=13$ :
$\dfrac{4}{10}=0{,}4$ et $\dfrac{6}{13}\approx0{,}46$ : quotients différents, les
droites ne sont pas parallèles. \\ \hline

\moment{7) Réinvestissement / travaux dirigés (5 min)} &
Le professeur propose l'exercice, fait passer un élève au tableau, puis corrige la rédaction. &
\textbf{Résolution} : les élèves cherchent, puis passent au tableau et rédigent la justification. &
\app \par
Dans chaque cas, les points sont alignés dans le même ordre. Les droites sont-elles
parallèles ? Justifier. \par
\textbf{1)} $AM=3{,}6$~cm ; $AB=9$~cm ; $AN=2{,}4$~cm ; $AC=6$~cm. \par
\textbf{Réponse :} $\dfrac{3{,}6}{9}=0{,}4$ et $\dfrac{2{,}4}{6}=0{,}4$ : quotients
égaux $\Rightarrow$ $(MN)\,//\,(BC)$ (réciproque de Thalès). \\ \hline

\moment{8) Remédiation éventuelle (3 min)} &
Le professeur reprend avec les élèves en difficulté  le calcul des quotients. &
Les élèves refont  un calcul de quotient simple. &
\titreR{Remédiation} : calculer $\dfrac{2{,}5}{10}$ et $\dfrac{3}{12}$ puis
comparer ; (\textit{réponse : $0{,}25=0{,}25$}). \\ \hline

\moment{9) Évaluation (formative) (2 min)} &
Le professeur interroge oralement, puis donne l'exercice de maison. &
Les élèves répondent à main levée, puis notent l'exercice de maison. &
\titreR{Évaluation formative} : citer les $4$ étapes de la justification. \par
(\textit{Réponse : ordre ; quotients ; comparaison ; conclusion}.) \par
\maison \par
$M\in[AB]$, $N\in[AC]$, dans le même ordre : $AM=2{,}8$~cm ; $AB=7$~cm ;
$AN=3{,}6$~cm ; $AC=9$~cm. \par
Les droites $(MN)$ et $(BC)$ sont-elles parallèles ? Justifier. \par
\textit{(Réponse attendue : oui, $\dfrac{2{,}8}{7}=\dfrac{3{,}6}{9}=0{,}4$ ;
réciproque de Thalès.)} \\ \hline
\end{longtable}

\begin{center}
\underline{\textcolor{rougetitre}{\large\bfseries Institutionalisation \; : \; (trace écrite de la leçon par le professeur)}}
\end{center}

% ═══════════════════════════════════════════════════════
%  SÉANCE 7 — SYNTHESE 7 + IV) PROGRAMME DE CONSTRUCTION
% ═══════════════════════════════════════════════════════
\begin{longtable}{|p{2.35cm}|p{5.35cm}|p{4.55cm}|p{13.85cm}|}
\hline
\seance{SEANCE 7} \newline \moment{1) Remobilisation des prérequis ou évaluation diagnostique (10 min)} &
Le professeur fait corriger l'exercice de maison au tableau, puis fait le point sur les erreurs. &
\textbf{Résolution} : deux élèves traitent l'exercice au tableau ; les autres corrigent. &
\titreR{Correction de l'exercice de maison} \par
$\dfrac{AM}{AB}=\dfrac{2{,}8}{7}=0{,}4$ et $\dfrac{AN}{AC}=\dfrac{3{,}6}{9}=0{,}4$ :
quotients égaux, points dans le même ordre $\Rightarrow$ $(MN)\,//\,(BC)$ d'après la
réciproque de Thalès. \\ \hline

\moment{2) Présentation de la situation (5 min)} &
Le professeur rappelle la situation  et lit l'énoncé de la Synthèse 7. &
Les élèves récapitulent la situation et notent l'énoncé. &
\titreR{Synthèse 7 (S7)} : sur le plan de l'entrée, on veut une nouvelle
configuration vérifiant $\dfrac{AM}{AB}=\dfrac{1}{3}$. \par
\textbf{1)} Écrire le programme de construction. \textbf{2)} Réaliser la figure et
justifier que la configuration obtenue a bien la propriété demandée. \\ \hline

\moment{3) Appropriation de la situation, compréhension de la tâche et de
l'organisation du travail (5 min)} &
Le professeur pose des questions de compréhension : « Quels instruments utiliser ? Dans quel ordre construire ? Comment justifier ensuite ? » Il organise ensuite les groupes. &
Les élèves répondent, clarifient la tâche et s'organisent en groupes. &
\titreR{Consignes} : ordonner les instructions (triangle, placement de $M$,
parallèle), construire, puis justifier avec la propriété de Thalès. \\ \hline

\moment{4) Résolution du problème : individuellement (5 min) puis en groupes (10 min)} &
Le professeur vérifie les productions, circule parmi les groupes et les guide sans donner la solution. &
\textbf{Résolution} : les élèves écrivent le programme, puis construisent d'abord individuellement et ensuite en groupes. &
\titreR{Synthèse 7 (S7) — Résolution} \par
\textbf{1)} tracer un triangle $ABC$ ; mesurer $AB$ ; placer $M$ sur $[AB]$ tel que
$AM=\dfrac{1}{3}AB$ ; tracer la parallèle à $(BC)$ passant par $M$ ; elle coupe
$[AC]$ en $N$. \par
\textbf{2)} $(MN)\,//\,(BC)$, $M\in[AB]$, $N\in[AC]$ : d'après la propriété de
Thalès, $\dfrac{AN}{AC}=\dfrac{AM}{AB}=\dfrac{1}{3}$ : la configuration a bien la
propriété demandée. \\ \hline

\moment{5) Synthèse et bilan du travail (3 min)} &
Le professeur fait reformuler la démarche, corrige puis valide. &
Les élèves reformulent  la démarche (programme, construction, justification). &
\titreR{Synthèse et bilan} : un programme de construction est une suite ordonnée
d'instructions, vérifiée par les propriétés du cours. \\ \hline

\moment{6) Institutionnalisation : trace écrite de la leçon par le professeur (7 min)} &
Le professeur présente la notion de programme de construction d'une configuration ; il fait écrire l'exemple au tableau, fait réaliser la construction, puis fait recopier la leçon. &
Les élèves écrivent le programme, réalisent la construction aux instruments, puis recopient la leçon. &
\titreR{Trace écrite} \par
\textbf{IV) Utilisation des définitions et propriétés pour justifier} \par
\textbf{1) Propriété d'une configuration — Programme de construction} \par
— Un \mclef{programme de construction} est une suite ordonnée d'instructions
utilisant les instruments (règle, compas, équerre) pour obtenir une configuration
donnée. \par
— \mclef{Justifier une propriété d'une configuration :} on utilise les définitions
et les propriétés du cours (propriété de Thalès, sa réciproque, quotients égaux). \par
\prop{\textbf{\textcolor{rougetitre}{Exemple (programme de construction)}} \par
Construire une configuration de Thalès vérifiant $\dfrac{AM}{AB}=\dfrac{1}{3}$ : \par
\textbf{1)} tracer un triangle $ABC$ quelconque ; \par
\textbf{2)} placer $M$ sur $[AB]$ tel que $AM=\dfrac{1}{3}\times AB$ (mesurer $AB$,
calculer, reporter) ; \par
\textbf{3)} tracer la parallèle à $(BC)$ passant par $M$ ; elle coupe $[AC]$ en $N$ ; \par
\textbf{4)} justifier : $(MN)\,//\,(BC)$ et $M\in[AB]$, $N\in[AC]$, donc
$\dfrac{AM}{AB}=\dfrac{AN}{AC}=\dfrac{MN}{BC}$ (propriété de Thalès) ;
ainsi $\dfrac{AN}{AC}=\dfrac{1}{3}$ : la configuration a bien la propriété demandée.} \\ \hline

\moment{7) Réinvestissement / travaux dirigés (5 min)} &
Le professeur propose l'exercice, fait écrire le programme, fait réaliser la construction, puis corrige. &
\textbf{Résolution} : les élèves écrivent le programme, construisent, puis passent au tableau. &
\app \par
\textbf{1)} Écrire le programme de construction d'une configuration dans laquelle
$\dfrac{AM}{AB}=\dfrac{3}{4}$ ($ABC$ triangle). \par
\textbf{Réponse :} tracer $ABC$ ; mesurer $AB$ ; placer $M$ tel que
$AM=\dfrac{3}{4}AB$ ; parallèle à $(BC)$ par $M$ ; elle coupe $[AC]$ en $N$ ; \par
justification : d'après la propriété de Thalès,
$\dfrac{AN}{AC}=\dfrac{MN}{BC}=\dfrac{3}{4}$. \\ \hline

\moment{8) Remédiation éventuelle (3 min)} &
Le professeur reprend avec les élèves en difficulté  l'ordre des instructions d'un programme. &
Les élèves remettent dans l'ordre les étapes d'un programme donné en désordre. &
\titreR{Remédiation} : remettre en ordre $4$ étapes mélangées d'un programme de
construction. \\ \hline

\moment{9) Évaluation (formative) (2 min)} &
Le professeur interroge oralement, puis donne l'exercice de maison. &
Les élèves répondent à main levée, puis notent l'exercice de maison. &
\titreR{Évaluation formative} : qu'est-ce qu'un programme de construction ? \par
(\textit{Réponse : une suite ordonnée d'instructions utilisant les instruments}.) \par
\maison \par
Écrire le programme de construction d'une configuration de Thalès vérifiant
$\dfrac{AM}{AB}=\dfrac{2}{5}$, puis la réaliser et justifier que $MN=0{,}4\times BC$. \par
\textit{(Réponse attendue : programme en $4$ étapes ; justification par la propriété
de Thalès.)} \\ \hline
\end{longtable}

\begin{center}
\underline{\textcolor{rougetitre}{\large\bfseries Institutionalisation \; : \; (trace écrite de la leçon par le professeur)}}
\end{center}

% ═══════════════════════════════════════════════════════
%  SÉANCE 8 — SYNTHESE 8 + RÉCAPITULATIF DE LA LEÇON
% ═══════════════════════════════════════════════════════
\begin{longtable}{|p{2.35cm}|p{5.35cm}|p{4.55cm}|p{13.85cm}|}
\hline
\seance{SEANCE 8} \newline \moment{1) Remobilisation des prérequis ou évaluation diagnostique (10 min)} &
Le professeur fait corriger l'exercice de maison au tableau, puis fait le point sur les erreurs. &
\textbf{Résolution} : deux élèves présentent leur programme et leur figure ; les autres corrigent. &
\titreR{Correction de l'exercice de maison} \par
Programme : tracer $ABC$ ; mesurer $AB$ ; placer $M$ tel que
$AM=\dfrac{2}{5}AB$ ; parallèle à $(BC)$ par $M$, elle coupe $[AC]$ en $N$. \par
Justification : propriété de Thalès $\Rightarrow \dfrac{MN}{BC}=\dfrac{AM}{AB}
=\dfrac{2}{5}=0{,}4$, donc $MN=0{,}4\times BC$. \\ \hline

\moment{2) Présentation de la situation (5 min)} &
Le professeur rappelle la situation  complète et lit l'énoncé de la Synthèse 8. &
Les élèves récapitulent la situation et notent l'énoncé. &
\titreR{Synthèse 8 (S8) : question de synthèse} \par
Le contrôleur mesure $BC=6$~m et la corde $MN=4$~m. \par
\textbf{1)} Vérifier par le calcul que $\dfrac{MN}{BC}=\dfrac{AM}{AB}
=\dfrac{AN}{AC}$. \par
\textbf{2)} En déduire la position de $(MN)$ par rapport à $(BC)$. \par
\textbf{3) } Que se passerait-il si la corde mesurait $MN=5$~m ? \\ \hline

\moment{3) Appropriation de la situation, compréhension de la tâche et de
l'organisation du travail (5 min)} &
Le professeur pose des questions de compréhension, puis organise les groupes pour cette question de synthèse. &
Les élèves répondent, clarifient la tâche et s'organisent en groupes. &
\titreR{Consignes} : utiliser la propriété de Thalès (1), la réciproque (2), puis
argumenter sur le cas $MN=5$~m (3). \\ \hline

\moment{4) Résolution du problème : individuellement (5 min) puis en groupes (10 min)} &
Le professeur vérifie les productions, circule parmi les groupes et les guide sans donner la solution. &
\textbf{Résolution} : les élèves cherchent d'abord individuellement, puis en groupes ; ils exposent leurs démarches. &
\titreR{Synthèse 8 (S8) — Résolution} \par
\textbf{1)} $\dfrac{MN}{BC}=\dfrac{4}{6}=\dfrac{2}{3}$ ;
$\dfrac{AM}{AB}=\dfrac{6}{9}=\dfrac{2}{3}$ ; $\dfrac{AN}{AC}
=\dfrac{5}{7{,}5}=\dfrac{2}{3}$ : vérifié. \par
\textbf{2)} quotients égaux et points dans le même ordre : $(MN)\,//\,(BC)$
(réciproque de Thalès). \par
\textbf{3)} $\dfrac{5}{6}\neq\dfrac{2}{3}$ : la corde ne serait plus parallèle à
$(BC)$ (quotients différents). \\ \hline

\moment{5) Synthèse et bilan du travail (3 min)} &
Le professeur organise le bilan de toute la leçon à partir des huit Synthèses. &
Les élèves récapitulent oralement les outils de la leçon. &
\titreR{Synthèse et bilan} : les Synthèses 1 à 8 ont permis d'identifier les
configurations (S1, S2), de calculer des grandeurs (S3, S4, S5) et de justifier
(S6, S7, S8). \\ \hline

\moment{6) Institutionnalisation : trace écrite de la leçon par le professeur (7 min)} &
Le professeur fait récapituler les outils de la leçon, puis fait recopier le récapitulatif. &
Les élèves recopient  le récapitulatif de la leçon. &
\titreR{Trace écrite : récapitulatif} \par
\textbf{V) Récapitulatif de la leçon} \par
— \mclef{Configurations de Thalès :} cas particulier du triangle ; cas général
(papillon, sommet $O$). \par
— \mclef{Propriété de Thalès :} $(MN)\,//\,(BC)\Rightarrow\dfrac{AM}{AB}
=\dfrac{AN}{AC}=\dfrac{MN}{BC}$ $\left(\right.$ou $\dfrac{OA}{OB}=\dfrac{OA'}{OB'}\left.\right)$
$\to$ calculs de longueurs, partage d'un segment selon une proportion. \par
— \mclef{Propriété réciproque :} même ordre $+$ quotients égaux $\Rightarrow$
parallélisme $\to$ justification de deux droites parallèles. \par
— \mclef{Programme de construction :} suite ordonnée d'instructions, vérifiée par
les propriétés du cours. \\ \hline

\moment{7) Réinvestissement / travaux dirigés (5 min)} &
Le professeur propose la situation, fait chercher les élèves, fait exposer quelques groupes, puis corrige. &
\textbf{Résolution} : les élèves cherchent, exposent leur démarche, puis se corrigent entre eux. &
\app \par
\textit{Un muret vertical de $1{,}2$~m projette une ombre de $1{,}8$~m ; au même
moment, un arbre projette une ombre de $7{,}5$~m (rayons du soleil parallèles).} \par
\textbf{1)} Identifier la configuration de Thalès. \quad
\textbf{2)} Calculer la hauteur de l'arbre. \par
\textbf{Réponses :} \textbf{1)} configuration générale (rayons parallèles). \par
\textbf{2)} $h=\dfrac{7{,}5\times1{,}2}{1{,}8}=5$~m. \\ \hline

\moment{8) Remédiation éventuelle (3 min)} &
Le professeur reprend avec les élèves en difficulté  le choix de l'outil adapté à la question posée. &
Les élèves refont le tri : pour calculer une longueur, on utilise la propriété de Thalès ; pour prouver un parallélisme, on utilise sa réciproque. &
\titreR{Remédiation} : associer chaque question (calculer une longueur / prouver un
parallélisme) à l'outil adapté. \\ \hline

\moment{9) Évaluation (formative) (2 min)} &
Le professeur pose une dernière question orale, donne l'exercice bilan, annonce l'évaluation prévue, puis conclut la leçon. &
Les élèves répondent, puis notent l'exercice de maison. &
\titreR{Évaluation formative} : citer les deux énoncés (propriété et réciproque). \par
\maison \par
\textbf{1)} $M\in[AB]$, $N\in[AC]$ : $AM=6$~cm ; $MB=4$~cm ; $AN=3{,}6$~cm ;
$BC=8$~cm. Calculer $AB$, $AC$ et $MN$ ($2$ décimales si besoin). \par
\textbf{2)} En déduire, à l'aide de la réciproque, si les droites $(MN)$ et $(BC)$ sont parallèles. \par
\textit{(Réponse attendue : $AB=10$~cm, $AC=6$~cm, $MN=4{,}8$~cm ; $\dfrac{AN}{AC}=0{,}6=\dfrac{AM}{AB}$, donc $(MN)\,//\,(BC)$.)} \\ \hline
\end{longtable}

\end{document}
